\section{Reproducing the PoC}
\label{app:reproduce}

\paragraph{Repository and commit.} \href{https://github.com/AdamKrysztopa/edu_proj}{\nolinkurl{github.com/AdamKrysztopa/edu_proj}}.
This report was built against commit \code{754e2e6} on branch \code{main}. All committed artifacts
cited in this report (JSON ledgers, gap maps, test suites) are reachable from that commit; later
commits may add report sections without changing any N2/N3 artifact discussed here.

\paragraph{Environment.} Each package (\code{residual}, \code{reconstruct}, \code{gapmap},
\code{instrument}) is a separate \code{uv} project (Python \(\geq\)3.12, its own
\code{pyproject.toml}), run with \code{uv run --directory <package> <command>}. No paid API key is
needed for anything below; \code{gapmap}'s default mode makes no network call at all, and
\code{reconstruct}'s and \code{instrument}'s test suites never call a live model
(\cref{app:claims}, rows E11--E14).

\paragraph{The four test commands.}
\begin{verbatim}
uv run --directory residual pytest -q      # 268 passed, ~1 s
uv run --directory reconstruct pytest -q   # 416 passed, ~1 s
uv run --directory gapmap pytest -q        # 172 passed, ~2 min
uv run --directory instrument pytest -q    # 311 passed, ~7 s
\end{verbatim}
All four were run live for this report on 2026-09-29 (\cref{app:claims}, rows E11--E14); the
counts above are the ones the report-wide consistency ruling requires, and they matched the
independently-run numbers exactly.

\paragraph{Package internals.} Module counts below include \code{__init__.py} for all three
packages. \repo{residual/}'s nine modules (1{,}281 lines): \code{vocab.py}
(closed enumerations), \code{provenance.py} (\code{Record}, \code{Agent}, \code{Source},
\code{Selector}, \code{Verdict}, \code{Verification}, \code{Evidence}, \code{Generation},
\code{Search}, \code{Derivation}), \code{claims.py} (\code{ClaimRecord}, \code{Scope},
\code{ResponseDistribution}, and the \code{label} property), \code{ledger.py} (\code{Ledger},
\code{Contradiction}, \code{Area}), \code{gates.py} (\code{evidential_gate}, \code{Measurement},
\code{measure()}, \code{Threshold}), \code{gapmap.py} (the N1 feature definitions that N2 computes
and N3 predicts from), \code{freeze.py}, and \code{residual.py} (\code{ResidualAccount},
\code{EstimatedResidual}). \repo{reconstruct/} has ten modules (4{,}803 lines), built around
\code{run.py} (the \code{reconstruct()} entry point) and \code{evidence.py} (independence
clustering); its third-party dependencies are \code{httpx}, \code{openai}, \code{tldextract}, and
\code{pypdf}. \repo{gapmap/} has twelve modules (2{,}986 lines), built around \code{lenses.py}
(the seven firing rules), \code{semantic.py} (closure judging), \code{checks.py} (the built-in
adversarial checks), \code{rank.py} (the breadth-rubric ranking and the display cap), and
\code{record.py} (the three-tier \code{Record} and \code{categorize}); it depends on nothing beyond
stdlib, \code{pydantic}, and \code{residual} (\cref{sec:architecture}).

\paragraph{Residual accounting and the frozen N1 schema, in full.} What exists in code:
\code{AreaFeatures} (per-area counts of labels, source kinds, tacitness, and contradictions,
\repo{residual/src/residual/gapmap.py}) and \code{ResidualAccount}/\code{EstimatedResidual}
(\repo{residual/src/residual/residual.py}), which name the fields an estimate would need --
\code{EstimatedResidual} allows \code{chao1}, \code{jackknife}, \code{mh}, and
\code{lincoln_petersen} as its declared estimator methods. What does not exist: any estimator that
actually runs one of those methods, and any World~C data to run it against
(\cref{sec:evidence-model}). Decision 0006 hash-freezes the schema, the gap-map feature set, and
the label-deciding code together in \repo{residual/frozen/n1.json}, which holds three SHA-256
hashes: \code{schema_sha256} (the JSON schema of \code{ClaimRecord}, \code{Ledger},
\code{Measurement}, and every vocabulary's enum values, stripped of prose), \code{features_sha256}
(the \code{AreaFeatures} field-name tuple), and \code{semantics_sha256} (a byte hash of
\code{vocab.py}, \code{provenance.py}, \code{claims.py}, \code{ledger.py}, \code{gates.py}, and
\code{residual.py} -- the code that decides labels, gates, and coverage). A plain \code{pytest}
architecture test fails if the live fingerprint diverges from the committed file. Changing any of
it needs a deliberate re-freeze (\code{python -m residual.freeze --write}) and a commit explaining
why.

\paragraph{Replaying an N3 gap map offline.} Every committed gap map replays byte-for-byte from
its ledger plus the committed closure-judge cache (\code{closure_judgements.json} in each
\repo{research/n3/} domain folder); no network call happens unless \code{--judge ollama} is added.
Pointing \code{--out} at a fresh, empty directory fails with \code{JudgeCacheMiss}, because the
cache path the CLI reads is \code{<out>/closure_judgements.json}: the cache must be copied into the
output directory first, and byte identity needs the exact domain strings and sidecar/sibling flags
below (verified by running this sequence):
\begin{verbatim}
OUT=/tmp/gapmap-replay; R=reconstruct/runs
for d in plc gdpr_v1 gdpr_v2; do
  mkdir -p $OUT/$d
  cp research/n3/$d/closure_judgements.json $OUT/$d/
done
cd gapmap
uv run python -m gapmap \
  --ledger ../$R/20260928T133336Z-59ec876b3d7e/ledger.json \
  --sidecar ../$R/20260928T133336Z-59ec876b3d7e/sidecar.json \
  --domain "PLC intermittent-fault diagnosis" --out $OUT/plc
uv run python -m gapmap \
  --ledger ../$R/20260928T135354Z-eb840a85b7d1/ledger.json \
  --sibling ../$R/20260928T145943Z-bc1cc9a30a02/ledger.json \
  --sidecar ../$R/20260928T135354Z-eb840a85b7d1/sidecar.json \
  --domain "GDPR DPIA (run v1)" --out $OUT/gdpr_v1
uv run python -m gapmap \
  --ledger ../$R/20260928T145943Z-bc1cc9a30a02/ledger.json \
  --sibling ../$R/20260928T135354Z-eb840a85b7d1/ledger.json \
  --sidecar ../$R/20260928T145943Z-bc1cc9a30a02/sidecar.json \
  --domain "GDPR DPIA (run v2)" --out $OUT/gdpr_v2
cd ..
for d in plc gdpr_v1 gdpr_v2; do
  cmp $OUT/$d/gapmap.json research/n3/$d/gapmap.json
  cmp $OUT/$d/gapmap.md   research/n3/$d/gapmap.md
done
\end{verbatim}
This reads the \emph{hashed} lens configuration (\code{config.py}), not a frozen one: the lexicons
and thresholds were tuned in-sample on all three ledgers together, and are hash-recorded
(\code{config_sha256}, below) so that a future change is visible, not blocked (\cref{sec:n3}).
Adding \code{--judge ollama} re-judges live against a local Ollama server running
\code{qwen2.5:7b-instruct}, which this report never does and which is not needed to check any
number in it. Separately, \code{uv run --directory reconstruct python -m reconstruct.verify_run
runs/<run>} re-checks every evidence item's span against the run's local snapshot; it needs the
local, git-ignored snapshots (below), so it cannot run from a fresh clone.

\paragraph{Headline numbers and the release.} Every headline number in both reports is listed, with
its artifact, field and file hash, in \repo{research/results_manifest.json}.
\code{python3 scripts/results_manifest.py --check} recomputes each value from the artifacts and fails on
any mismatch. Repository links in both reports point at the audited commit named on the title page,
not at the moving \code{main} branch; an earlier N3 README that still quoted a superseded
lexical-closure result showed why.

\paragraph{The Jev spike.} Its analyses re-run offline from the committed cache
(\repo{research/spikes/jev/cache}), with no provider call for any cached request; see
\repo{research/spikes/jev/RESULTS.md} for the subcommands.

\paragraph{The N2 ledger run IDs used in this report.}
\begin{itemize}
  \item PLC: \code{reconstruct/runs/20260928T133336Z-59ec876b3d7e}
  \item GDPR v1: \code{reconstruct/runs/20260928T135354Z-eb840a85b7d1}
  \item GDPR v2: \code{reconstruct/runs/20260928T145943Z-bc1cc9a30a02} (defect-compromised per
    \cref{app:claims} row E23; not N3-admissible on its own, used only for the cross-run-stability
    check, row E34)
\end{itemize}

\paragraph{Ledger and config hashes (from the committed gapmap.json files).}
Every gap map records the SHA-256 of the exact ledger it was built from and of the exact lens
configuration (\code{config.py}) that produced it, so a mismatch between a recomputed hash and the
one below means the input has changed since this report was written.

\begin{center}
\footnotesize
\begin{tabular}{@{}l p{5.4cm} p{5.4cm}@{}}
  \toprule
  \textbf{Domain} & \textbf{ledger\_sha256} & \textbf{config\_sha256} \\
  \midrule
  PLC & \nolinkurl{3d724d8ddfb78f43240f853e7dee8702cd78d72b1d0856e9f1b70be7f90e49bd} &
    \nolinkurl{35fd0af25618f8d6e0d1c2c3707d7e932afe2ce72dbd01e956034b43bfeb04c1} \\
  GDPR v1 & \nolinkurl{cdd8e42e0a140e5d75cd9e83178479fa61ba55ed3aa0a701731f0ea7e68f004b} &
    \nolinkurl{35fd0af25618f8d6e0d1c2c3707d7e932afe2ce72dbd01e956034b43bfeb04c1} \\
  GDPR v2 & \nolinkurl{1a1d1ca74b82a933496456b0bbb2b957568453016ce160c1f57c3786dcc6788b} &
    \nolinkurl{35fd0af25618f8d6e0d1c2c3707d7e932afe2ce72dbd01e956034b43bfeb04c1} \\
  \bottomrule
\end{tabular}
\end{center}

All three domains share one \code{config_sha256}: the same \emph{hashed, not yet frozen}, lens
configuration -- tuned in-sample on these same three ledgers jointly -- was applied to all three,
so any difference in their gap maps comes from the ledgers, not from per-domain tuning of the
lenses. Hashing makes a future change visible; it does not yet block one the way decision 0006's
freeze blocks an undeclared change to the N1 schema (\cref{sec:evidence-model}).

\paragraph{What cannot be reproduced offline.}
\begin{itemize}
  \item \textbf{Any live model call.} E-LIVE, E-PLANT, E-ABST and a \code{--judge ollama} run all
    need a paid model key (OpenRouter) or a local Ollama server; none of this report's numbers
    depend on running one, and re-running E-LIVE would not reproduce the same run in any case,
    since it fetches the live web (\cref{sec:threats}).
  \item \textbf{Re-slicing evidence spans on a fresh clone.} \code{.gitignore} excludes
    {\small\nolinkurl{reconstruct/runs/**/snapshots/}}: the raw fetched documents behind each
    \code{evidence[].selector.locator} are not in the repository, only their SHA-256 hash. On the
    machine that ran E-LIVE, those snapshots are still present (git-ignored, not deleted), so
    \code{reconstruct.verify_run} re-slices all 909 evidence items exactly, and the closeout's
    71-of-71 sample (\cref{app:claims}, row E25) checked the same way. From a fresh clone,
    with no snapshots, a byte-offset \code{locator} (e.g.\
    \code{sha256:14b9bb35...;char=9061,9169}, \cref{app:records}) can be checked only against a new
    fetch of the same URL, and only if the page has not changed since 2026-09-28.
  \item \textbf{The E-PLANT/E-ABST sealed answer keys and gated-arm data.} These live under
    \code{.private/} and \code{sessions/}-style paths that are deliberately excluded from every
    tool in this project (\repo{CLAUDE.md}, \enquote{Off limits}), so this report's own authors
    could not, and a reader cannot, re-open them; only the aggregate numbers already committed to
    \repo{research/n2/e_plant_eabst_report.md} are available.
\end{itemize}
